% !TeX root = ../papers/nc-p3.tex
\section{曲線状有限窓での母線の排除と三重点の取り直し}
\label{sec:finite-curvilinear-pencils}
\begin{proposition}\label{prop:curvilinear-marked-pencils}
全中央支持4の有限窓の全局所環を曲線状とし、単純印の二つの四空間が五空間 $E$ を張るとする。
$\mathbf P(E)$ の階数一切断は母線を含まない。切断が被約六点なら、
印と異なる非単純点の少なくとも一つで $A,B$ の剰余値は比例する。
切断が平面二次曲線と二つの孤立被約点の非交和なら偶側対称である。
\end{proposition}
\begin{proof}[付値で測った鉛筆]
前節と同じく $h$ を $E_B$ の式とする。$h(A)$ の零化イデアルで定まる残余 $Z$ は長さ13以上。
印を除いた局所環 $R=k[t]/t^m$ で、$A,B$ が剰余で独立なら元の鉛筆を $R$ 上で使う。
比例なら、$A$ の単元成分に合わせて単元 $\lambda$ を選び、
$B=\lambda A+t^a V$ とする。差が零なら $h(A)=0$ で外積も零である。
そうでなければ $1\le a<m$ で、$A,V$ は $R'=k[t]/t^{m-a}$ 上で独立かつ
共通左右核を持つ。これは $A$ の選んだ成分で差が零、別の成分でちょうど位数 $a$ となるためである。
$h(A)=-t^a h(V)/\lambda$ より $Z$ の局所因子は $R'$ の商である。
独立の場合も $R'=R$ として、$Z$ 上には常に直和因子である二次元鉛筆がある。

母線 $ku\otimes T$ を仮定し、定理\ref{thm:finite-five-space-rank-one-curves}の
二行三列への射影 $E\to\bar E$ を用いる。$Z$ の一因子 $k[t]/t^s$ での位数は $s$ で打ち切る。
右評価 $\bar E\to k^2$ は剰余で非零である（零なら元の右評価階数は1以下）。
$\dim\bar E=2$ なら単元の軸を掃き出すと核は一本の自由直線に入る。
独立鉛筆の二本の射影を、その直線上で単元成分を持つ係数対で結合して零にできる。
二係数の一方が他方を割るからである。元の結合は単元成分を持ち $ku$ に値を取るので $l(u)=0$。
これだけでは印以外の長さ6までである。

$\dim\bar E=3$ では右評価の二次小行列の最小位数を $b$ とする。
一つの単元の軸の掃き出しと主イデアル性により行列を対角 $1,t^b$ と零列にできる。
鉛筆の二本の射影は第一成分零、第二成分が $(t^{s-b})$ にあり、
第三成分を上と同じく原始的な係数対で打ち消せる。元の結合を $t^{s-b}$ を法として取れば
$ku$ に値を持つ。従って
\[
                \operatorname{ord} l(u)+b\ \ge\ s .
\]
$\bar E$ の平面の階数一切断が有限なら、右評価の二次小行列の零点スキーム $Y$ も有限である。
実際、評価核が二次元となる各右方向は、平面の有限曲線状長3切断の
長さ2部分の割線で一意に決まる（同定理の母線の議論）。
$Y\subset\mathbf P^2$ の長さは3である。高さ2の極大小行列イデアルの
Hilbert--Burch分解は $0\to\mathcal O(-3)^2\to\mathcal O(-2)^3\to\mathcal I_Y\to0$ である。
空の場合も上界だけなら自明である。長さ $e_i$ の各局所因子のイデアルは
その点の極大イデアルの $e_i$ 乗を含む。従って全因子での $b$ の和は、
対応する右方向の逆像長を各 $e_i$ 倍した和以下である。
一定の右方向は印を通る直線の印以外の部分で長さ2以下ゆえ、合計は6以下である。
階数一切断が正次元なら、同じ平面の分類により $\bar E=\langle e_1v^t,e_2v^t,C\rangle$。
評価列は $v(r)e_1,v(r)e_2,Cr$、かつ $v(r)=0$ なら $Cr$ の剰余は非零なので、
小行列イデアルは層として $(v(r))$ である。この場合も印を通る平面で $\sum b\le6$ となる。
$\sum\operatorname{ord} l(u)\le6$ と合わせると $13\le\operatorname{length}Z\le12$ に矛盾する。

被約六点の場合、$R'$ 上の鉛筆の行列式で切る二点因子は定数六点の部分であり、鉛筆はその一定割線である。
$h(A)\ne0$ ならこの割線は $E_B$ 内にない。比例の剰余の場合も、
$A,V$ がともに $R'$ 上で $h$ を消せば $t^a h(V)=0$ に反する。
$B$ の $R'$ 上の像は従って割線と超平面の固定交点となり、閉埋込みより $R'$ は長さ1。
$Z$ の長さ13以上が互いに異なる割線を与え、前節の六点の応力の証明により一つ以上で
両超平面の交点が等しい。その点は単純になれない。

最後に二次曲線と被約二点の非交和では、$R'$ 上の鉛筆の二点因子が孤立点を含めば
その点は単根として一定に持ち上がる（被約孤立点に二重接線はない）。
その場合は二つの孤立点の包で割った商の外積が零となる。
残る鉛筆は相対的にも曲線の平面 $D$ 内にあり、左核は一定方向になる。
比例の剰余なら外積に $t^a$ が掛かるので、$R'$ 上で得たこの固定方向のイデアルによる
零化は元の環にも戻る。印で外積は零である。前節の商の三成分と長さ2の判定で偶側対称となる。
\end{proof}

\begin{theorem}\label{thm:finite-single-triple-five-space}
全中央支持4の有限窓で点スキームが十七単純点と一つの曲線状三重点からなるとき、
反対階数10なら少なくとも片側対称であり、反対階数11は不可能である。
\end{theorem}
\begin{proof}
単純印で五空間が出る場合を調べる。$J$ の外に少なくとも $18-7=11$ 個の異なる支持があり、
各々で元の二本が独立である。以下、対称性がまだ得られないと仮定する。

切断が有限非被約なら、接次元2の点では全接線がその接平面にある。
両超平面がこの平面を含めば $J$ の外には寄与しない。それ以外は片方で直線以下を切り
高々3点しか捉えず、他の割線の高々4本を加えても足りない。
全点曲線状なら11本に達するのは単純四点と一つの二重点の場合だけで、
十本の異なる二点の組と一つの接線が全て使われる。
各超平面内の節点が四つ以上なら、単一四空間の長さ4判定が対称性を与える。
一つから三つならその一つと外の異なる二節点との二割線は同じ点で超平面に交わり、
座標埋込みの単射性に反する。従って両超平面とも全節点を避ける。
$Z$ の別の支持を加える余地はない。実際そこでも先の残余鉛筆の剰余割線は既に使われ、
その $E_B$ との唯一の交点が別の支持と一致してしまう。
従って $Z$ の長さ13を得るには三重点が十一支持の一つであり、そこで全長3を使う。
しかし元の独立鉛筆の行列式因子は長さ3の基底上では存在してもその割線の動きは
高々長さ2の基底を経由する。異なる二節点の場合、各根は単根で一意に分離し、
動ける一根は長さ2の節点への写像だけで決まる。一節点に重なる場合は一定の
長さ2節点全体に入る有限平坦次数2の閉部分であり、同次数ゆえ等しい。
いずれでも $E_B$ との単一交点 $B$ の閉埋込みは長さ3を通せない。

正次元切断で整曲線が全体を張る場合と二つの二次曲線が張る場合は、既出の長さ4判定を使う。
命題で母線はない。平面二次曲線の模型で $Q=h=0$ なら全 $E$ が対称にでき、
$E_A$ の階数一部分は $\mathbf P^2$ の二次Veroneseに二超平面を課したものとなる。
これは有限である（単一四空間の正次元切断なら母線を含む）。よって二つの二次式の完全交叉で長さ4。
$Q\ne0$ の二つの孤立被約点の場合も命題で済む。残る場合、前節の支持上界では11点に達するには
孤立点一つと外向き接点一つが必要である。$D$ 内の鉛筆で2点、孤立点を通る鉛筆で6点、
外向き接線で3点の各上界を同時に達し、三集合は互いに素となる。
最後の三点は両超平面の少なくとも一方と固定接平面との交線にあり、
前の二群もそれぞれ印を通る直線・平面の印外の長さ上限を使い切っている。
$Z$ の他の支持は残余鉛筆によってこの三条件の一つに入らねばならず、加える余地がない。
再び三重点は十一支持の一つで、そこに元の独立鉛筆がある。

この場合の外向き接点での完備局所環を求める。二つの相異なる根の一つは $h=0$ の単根で、
$h$ は非零一次式である。接点の左上を1とし、階数一行列を
\[
 \begin{pmatrix}1&x&z\\ y&yx&yz\\ w&wx&wz\end{pmatrix}
\]
と書く。商座標に $z,h$ を取れる。実際 $h=0$ の方向で $c,d$ の第一成分だけが非零となるよう
曲線の基底を選べるからである。$w=bz+ch$, $yz=\alpha h$, $wx=\beta h$,
$x-y=kz+k'h$, $h=wz$ で $b\ne0$。$Q=h((\beta-\alpha b-k)z+(-\alpha c-k')h)$
より $\beta-\alpha b-k\ne0$ である。$w=z b/(1-cz)$ と代入し、
$z(x-y)$ を二通りに計算すれば $z^2$ と単元の積が零となる。
従って局所環はちょうど $k[[x,z]]/(xz,z^2)$ で、$D$ は $z=0$ である。
孤立点はもう一つの単根なので被約であり、曲線の残りでは接次元1なので滑らかである。

三重点での鉛筆が孤立点を通れば、単根の分離によって全長でその点を通り、六点で飽和した
平面に入ってしまう。外向き接線の場合、行列式の相対長さ2の代数は剰余で $1,z$ が基底なので
全体でもこれが基底である。$x=A_0+A_1z$ と書けば $xz=z^2=0$ より $A_0=0$。
従って長さ2因子は固定接平面にスキームとして入り、その包である鉛筆もそこに入る。
三点で飽和した交線に長さを加えることはできない。
$D$ 内の場合、剰余鉛筆が特別接点に集中した二重点を切れば、それは曲線の接線でやはり固定接平面にあり、
既にその交線上にある別の三点と両立しない。特別接点を避けるなら全長で $D$ 内にある。
特別接点ともう一つの相異なる点を切れば単根二つに分離し、特別接点に延びる根の
$z$ は $z^2=0$ より $k[t]/t^3$ の $(t^2)$ に入る。
従って少なくとも $t^2$ を法として二根とも $D$ 内にあり、鉛筆全体とその固定左核も
その長さ2で $D$ 内の条件を満たす。二つの異なる支持で飽和した印外の直線に更に一重を加える矛盾となる。

以上で各単純印の五空間は、対称性がなければ被約六節点の場合に限られ、
命題よりその割線の一致点は唯一の三重点である。階数10または11では各単純印で一方の和の次元が
5以下で、4の場合は既に対称性が従う。十七印の少なくとも九印で同じ側の一致が生じるから、
前節と同様に三重点に印を移すと共通四空間となる。曲線状印での共通切断判定が対称性を与える。
\end{proof}
